
Reinforcement learning from human feedback (RLHF) can alter which answers a language model selects for multiple-choice questions (MCQs), yet no method exists to identify alignment-vulnerable items before fine-tuning begins. This paper examines whether the pre-alignment confidence margin --- the gap between a base model's top-1 and top-2 log-probabilities for a given item --- predicts post-alignment argmax instability. Across three benchmarks (MMLU, TruthfulQA, and ARC-Challenge), the pre-alignment margin predicts post-alignment argmax flips for a DPO-aligned 7B model with area under the receiver operating characteristic curve (AUROC) of 0.867 (MMLU), 0.803 (TruthfulQA), and 0.909 (ARC-Challenge), without retraining between benchmarks. A mixed-effects logistic regression yields $\beta _1 = -4.33$ (p $\approx 10^{-227}, \eta^2$ = 0.289) for this pair. The flip rate is 12.5\% across 14,042 MMLU items. Additionally, alignment-induced logit perturbations are structurally non-isotropic: the dominant eigenvalue of the logit-delta covariance is 2.90 times the mean of the remaining eigenvalues for DPO alignment and 4.58 times for SFT alignment, compared to a ratio of approximately 1.13 for an isotropic Gaussian control. An unexpected finding is that DPO-aligned models exhibit a monotone quintile trend in logit delta variance --- rising from 0.707 in the lowest confidence quintile to 3.384 in the highest --- while SFT produces a flat profile (0.223 to 0.281). These findings are limited to MCQ settings, to two model pairs (one DPO, one SFT), and the originally planned PPO comparison could not be executed due to model unavailability. Two planned mechanistic analyses (H-M3, H-M4) were not completed.

